
\section{Frank--Wolfe \label{sec:FrankWolfe}}

Mirror descent is not suitable for all spaces. The guarantee of mirror
descent crucially depends on the fact that there is a $1$-strongly
convex mirror map $\Phi$ such that $\max_{x}\Phi(x)-\min_{x}\Phi(x)$
is small on the domain. For some domains such as $\{x:\|x\|_{\infty}\leq1\}$,
the range $\max_{x}\Phi(x)-\min_{x}\Phi(x)$ can be large.
\begin{lem}
Let $\Phi$ be a $1$-strongly convex function on $\Rn$ over $\|\cdot\|_{\infty}\leq1$.
Then, 
\[
\max_{\|x\|_{\infty}\leq1}\Phi(x)\geq\min_{\|x\|_{\infty}\leq1}\Phi(x)+\frac{n}{2}.
\]
\end{lem}

\begin{proof}
By the strong convexity of $\Phi$, we have that 
\begin{align*}
\E_{s_{k}\in\{\pm1\}\ \forall k\in[n]}\Phi(s_{1},\cdots,s_{n}) & \geq\E_{s_{k}\in\{\pm1\}\ \forall k\in[n-1]}\Phi(s_{1},\cdots,s_{n-1},0)+\E_{s_{n}}\left\langle \nabla\Phi(s_{1},\cdots,s_{n-1},0),s_{n}\right\rangle +\frac{1}{2}\\
 & =\E_{s_{k}\in\{\pm1\}\ \forall k\in[n-1]}\Phi(s_{1},\cdots,s_{n-1},0)+\frac{1}{2}\\
 & \vdots\\
 & \geq\Phi(0)+\frac{n}{2}.
\end{align*}
\end{proof}
\begin{rem*}
This inequality is tight because $\frac{1}{2}\|x\|^{2}$ is $1$-strongly
convex in $\|\cdot\|_{\infty}$ and its value is between $0$ and
$\frac{n}{2}$.
\end{rem*}

\subsection{Algorithm}

Now we give another geometry dependent algorithm that relies on a
different set of assumptions. The problem we study in this section
is of the form
\[
\min_{x\in\mathcal{D}}f(x)
\]
for $f$ such that $\nabla f$ is ``Lipschitz'' in a certain sense.
The algorithm reduces the problem to a sequence of linear optimization
problems.

\begin{algorithm2e}[H]

\caption{$\mathtt{FrankWolfe}$}

\SetAlgoLined

\textbf{Input:} Initial point $x^{(0)}\in\Rn$, step size $h>0$.

\For{$k=0,1,\cdots,T-1$}{

Compute $y^{(k)}=\arg\min_{y\in\mathcal{D}}\left\langle y,\nabla f(x^{(k)})\right\rangle $.

$x^{(k+1)}\leftarrow(1-h_{k})x^{(k)}+h_{k}y^{(k)}$ with $h_{k}=\frac{2}{k+2}$.

}

\Return $x^{(T)}$.

\end{algorithm2e}

\subsection*{Analysis}
\begin{thm}
Let $f$ be a convex function on a convex set $\mathcal{D}$ with
a constant $C_{f}$ such that
\[
f((1-h)x+hy))\leq f(x)+h\left\langle \nabla f(x),y-x\right\rangle +\frac{1}{2}C_{f}h^{2}.
\]
Then, for any $x,y\in\mathcal{D}$ and $h\in[0,1]$, we have
\[
f(x^{(k)})-f(x^{*})\leq\frac{2C_{f}}{k+2}.
\]
\end{thm}

\begin{rem*}
If $\nabla f$ is $L$-Lipschitz with respect to the norm $\|\cdot\|$
over the domain $\mathcal{D}$, then $C_{f}\leq L\cdot\text{diam}_{\|\cdot\|}(\mathcal{D})^{2}.$
\end{rem*}
\begin{proof}
By the definition of $C_{f}$, we have that
\[
f(x^{(k+1)})\leq f(x^{(k)})+h_{k}\left\langle \nabla f(x^{(k)}),y^{(k)}-x^{(k)}\right\rangle +\frac{1}{2}C_{f}h_{k}^{2}.
\]
Note that 
\[
f(x^{*})\geq f(x^{(k)})+\left\langle \nabla f(x^{(k)}),x^{*}-x^{(k)}\right\rangle \geq f(x^{(k)})+\left\langle \nabla f(x^{(k)}),y^{(k)}-x^{(k)}\right\rangle 
\]
where we used the fact that $y^{(k)}=\arg\min_{y\in\mathcal{D}}\left\langle y,\nabla f(x^{(k)})\right\rangle $
and that $x^{*}\in\mathcal{D}$. Hence, we have that
\[
f(x^{(k+1)})\leq f(x^{(k)})-h_{k}(f(x^{(k)})-f(x^{*}))+\frac{1}{2}C_{f}h_{k}^{2}.
\]
Let $e_{k}=f(x^{(k)})-f(x^{*})$. Then, 
\[
e_{k+1}\leq(1-h_{k})e_{k}+\frac{1}{2}C_{f}h_{k}^{2}.
\]
Note that $e_{0}=f(x^{(0)})-f^{*}\leq\frac{1}{2}C_{f}$. By induction,
we have that $e_{k}\leq\frac{2C_{f}}{k+2}$.
\end{proof}
\begin{rem*}
Note that this proof is in fact the same as Theorem \ref{thm:gd_convex}.

\end{rem*}


